\subsection{The remaining order-five strict group is cyclic}

Write $\chi_p$ for the $p$-adic cyclotomic character and
$\omega$ for its Teichm\"uller reduction.  For an even integer
$e\ge2$ set
\[
 D=(\mathbf Q_p/\mathbf Z_p)(\chi_p^{-e}),\qquad
 \mathcal S_p(e)=
 \ker\!\left(H^1(G_{\mathbf Q,\{p\}},D)
       \longrightarrow H^1(G_{\mathbf Q_p},D)\right).
 \tag{O5S.1}
\]
Thus the classes are unramified away from $p$ and are zero
on the full decomposition group at $p$.
The carried VK-LOCAL comparison identifies this with the
source's $\operatorname{ShaU}(\chi_p^{-e};2p)$: its additional
decomposition condition at $2$ is automatic for the divisible
coefficients in (O5S.1).

\begin{corollary}[KT-ORDER-FIVE-STRICT-CYCLIC]
Let $p>3$ be prime, let $e\ge2$ be even, and suppose
$\operatorname{ord}(\omega^e)=5$.
Put $L=\mathbf Q(\mu_p)$, $\Delta=\operatorname{Gal}(L/\mathbf Q)$,
$A_p=\operatorname{Cl}(L)\{p\}$ and $\alpha=\omega^{-e}$.
There is a natural injection
\[
 \mathcal S_p(e)[p]\lhook\joinrel\longrightarrow
 \operatorname{Hom}_{\mathbf F_p[\Delta]}
       \bigl(A_p/pA_p,\mathbf F_p(\alpha)\bigr).
 \tag{O5S.2}
\]
Consequently
\[
            \dim_{\mathbf F_p}\mathcal S_p(e)[p]\le1.
 \tag{O5S.3}
\]
Under the carried finiteness and full-group VK-PT comparison,
there is an integer $a\ge0$ such that
\[
 \mathcal S_p(e)\simeq\mathbf Z/p^a\mathbf Z,
 \qquad K_{2e}(\mathbf Z)\{p\}\simeq\mathbf Z/p^a\mathbf Z.
 \tag{O5S.4}
\]
The argument asserts neither $a=0$ nor an upper bound on $a$.
There is no restriction $e<p$.
\end{corollary}

\begin{proof}
Let $M=\mathbf F_p(\alpha)$ and identify it with $D[p]$.
Since $\alpha$ has order five, the cyclotomic image of local
inertia at $p$ contains an element $g$ for which
$\chi_p^{-e}(g)-1\in\mathbf Z_p^\times$.
It follows that
\[
 H^0(G_{\mathbf Q},D)=H^0(G_{\mathbf Q_p},D)=0.
 \tag{O5S.5}
\]
The exact sequence
\[
          0\longrightarrow M\longrightarrow D
            \xrightarrow{\ p\ }D\longrightarrow0
 \tag{O5S.6}
\]
therefore identifies $H^1(G_{\mathbf Q},M)$ with
$H^1(G_{\mathbf Q},D)[p]$, and similarly gives an injection
$H^1(G_{\mathbf Q_p},M)\hookrightarrow
H^1(G_{\mathbf Q_p},D)$.
In particular, the finite-coefficient class corresponding to
$s\in\mathcal S_p(e)[p]$ is zero on $G_{\mathbf Q_p}$.
For any prime $q\ne p$, inertia acts trivially on $D$,
and $H^0(I_q,D)=D$ is $p$-divisible.
The same exact sequence gives an injection
\[
                    H^1(I_q,M)\hookrightarrow H^1(I_q,D).
 \tag{O5S.7}
\]
The finite-coefficient class is thus unramified at every
$q\ne p$ as well.  These arguments also justify using the
unramified-outside-$p$ quotient in (O5S.1).

The coefficient $M$ becomes trivial on $G_L$.
Since $|\Delta|=p-1$ is prime to $p$, inflation--restriction
and averaging give an injective restriction
\[
 H^1(G_{\mathbf Q},M)\lhook\joinrel\longrightarrow
 H^1(G_L,M)^\Delta.
 \tag{O5S.8}
\]
The restricted class is a continuous homomorphism $G_L\to M$.
It is unramified at every finite place of $L$: at the places
above $p$ this follows from the full local vanishing already
proved, and elsewhere from (O5S.7).
Global class field theory therefore makes it factor through
$A_p/pA_p$.
The equivariance convention is explicit: for a lift of
$\delta\in\Delta$ and $g\in G_L$, a restricted cocycle $c$
satisfies
\[
                  c(\delta g\delta^{-1})=\alpha(\delta)c(g).
 \tag{O5S.9}
\]
Functoriality of the Artin map gives precisely the Hom group
in (O5S.2), with character $\alpha$, not its inverse.
This proves the asserted injection.

Integral Teichm\"uller idempotents identify the dimension of
the Hom group with
\[
          \dim_{\mathbf F_p}
             A_p^{[-e]}/pA_p^{[-e]}.
 \tag{O5S.10}
\]
The character $\alpha$ is one of the four order-five characters
in KT-ORDER-FIVE-LENGTH-ONE.  That proposition bounds the sum
of their $\mathbf Z_p$-lengths by one, and proves (O5S.3).

The carried VK-PT theorem is a full group isomorphism
\[
              \bigl(K_{2e}(\mathbf Z)\{p\}\bigr)^\vee
                       \simeq\mathcal S_p(e),
 \tag{O5S.11}
\]
not merely a comparison of vanishing assertions.
The positive even-degree $K$-group here is finite by the
carried finiteness input.  Hence $\mathcal S_p(e)$ is a finite
abelian $p$-group.  For any such group $B$, the dimensions of
$B[p]$ and $B/pB$ are equal; dimension at most one means that
$B$ is cyclic.  Its Pontryagin dual is cyclic of the same
order.  This proves (O5S.4).
\end{proof}

\paragraph{Why the class-group length does not give the same exponent.}
The proof uses the fact that $D[p]$ is trivial on $G_L$.
This does not hold for all higher torsion chunks of $D$.
More precisely, put $r=v_p(e)$.  The image of $G_L$ under
$\chi_p$ is $1+p\mathbf Z_p$, and the lifting-the-exponent
identity gives
\[
                  v_p((1+p)^e-1)=r+1.
 \tag{O5S.12}
\]
Thus $D[p^{r+2}]$ has nontrivial $G_L$-action; when $p\nmid e$
this already happens for $D[p^2]$.
There is consequently no all-level replacement of (O5S.2)
by homomorphisms from the same ordinary ideal class group.
The finite groups $\mathbf Z/p^a\mathbf Z$ all have
one-dimensional $p$-torsion when $a\ge1$.
Therefore (O5S.3) alone does not bound the integer $a$.
This is the precise limitation of the comparison proved here;
no claim that an additional argument cannot bound $a$ is made.

\paragraph{The ordinary $R=T$ cyclicity hypothesis remains distinct.}
Suppose, in addition, that the hypotheses of
KT-RESIDUAL-CYCLICITY-SCOPE hold: the prime splits in
$E=\mathbf Q(\sqrt{-3})$, and the specified ordinary
irreducible rank-two congruence supplies its Urban class.
With the notation of that theorem, one has
\[
 \dim_k\mathcal H_E^{c=-1}\ge1,
 \qquad
 \dim_k\mathcal H_E^{c=+1}
        \ge\dim_{\mathbf F_p}\mathcal S_p(e)[p].
 \tag{O5S.13}
\]
If the rational group in (O5S.4) is nonzero, its $p$-torsion
has dimension exactly one, and the two eigenspaces in
(O5S.13) give $\dim_k\mathcal H_E\ge2$.
Thus the new one-generator conclusion for the rational
group does not supply the one-dimensional \emph{total}
ordinary Selmer hypothesis in the inspected Berger--Klosin
comparison.  Outside those Urban-branch hypotheses,
(O5S.3) still gives no upper bound for the additional
$c=-1$ eigenspace.  The previously checked weight-range,
ordered-flag, and divisible-Selmer-bound hypotheses of
those comparisons are unchanged.  No $R=T$ theorem is
being applied in this corollary.

\paragraph{Exact carried comparison and source scope.}
The full-group comparison (O5S.11) and the equality of
the local conditions are the accepted K17 theorem
\emph{VK-PT and VK-LOCAL}, equations (K5)--(K7), in
\texttt{RETURN\_K17\_checkpoint\_08\_20261121T0616Z/K17.tex},
lines 788--898.  Its entire proof was reread for this
application; its named foundational inputs remain carried.
The $\operatorname{ShaU}$ definition and its exact
finite-coefficient local transport were also recorded in
KT-ODD-CONORMAL-SELMER-EQUIVALENCE.
The order-five class bound is the separately checked
KT-ORDER-FIVE-LENGTH-ONE.  The additional ordinary
eigenspace in (O5S.13) is used only with all the hypotheses
of KT-RESIDUAL-CYCLICITY-SCOPE, not asserted for an arbitrary
prime or an arbitrary order-five character.
